% ----------------------------------------------------------
\chapter{Desenvolvimento}\label{cap:desenvolvimento}
% ----------------------------------------------------------

\section{Ambiente e Tecnologias}\label{sec:ambiente}

\section{Arquitetura do Software}\label{sec:arquitetura}

\section{Reprodutibilidade dos Resultados}\label{sec:reprodutibilidade}

\section{Testes Automatizados}\label{sec:testes}

\section{Rastreabilidade das Decisões}\label{sec:rastreabilidade}

\section{Painel de Resultados}\label{sec:painel}

\section{Implantação}\label{sec:implantacao}

O modelo final de melhor desempenho será encapsulado em um protótipo funcional por meio do \textit{framework} FastAPI \cite{ramirez2018fastapi}, cuja função é expor os \textit{endpoints} de classificação individual a sistemas externos, viabilizando a futura integração com bases de dados consolidadas da APS brasileira, a exemplo do Prontuário Eletrônico do Cidadão (PEC) no âmbito do SISAB \cite{brasil_esus_aps}. Adicionalmente, será desenvolvida uma interface gráfica por meio do \textit{framework} Streamlit, cuja função é proporcionar uma camada de interação visual para fins de demonstração técnica e validação acadêmica do protótipo, tornando acessível a sua utilização. Cabe ressaltar que o foco arquitetural do presente trabalho reside no desenvolvimento do motor classificatório (\textit{back-end}) e da camada de explicabilidade, sendo a interface gráfica apresentada como componente complementar de validação operacional, e não como ferramenta clínica validada ou sistema operacionalmente integrado à rede pública de saúde.
